\documentclass{article}
\usepackage[legalpaper, margin=1in]{geometry}

\title{Alpha Sigma Phi Beta Psi Constitution}
\author{Prudential Board of Alpha Sigma Phi Beta Psi}
\date{January 2026}

\begin{document}

\maketitle

\section*{Preamble}
\hspace*{1cm}We the brothers of the Beta Psi Chapter of Alpha Sigma Phi, in order to foster education, maintain charity, promote patriotism,  scholarship, assist in the building of character to promote college loyalties, perpetuate friendships, cement social ties within its membership, develop independence, do ordain and establish this Constitution for this chapter of the fraternity.

\section*{Article 1 - Name}
\hspace*{1cm}This chapter of Alpha Sigma Phi has been designated the Beta Psi Chapter, and is an active chapter as defined by and under the authority of Alpha Sigma Phi Fraternity.

\section*{Article 2 - Membership}
\hspace*{1cm}Membership into the Beta Psi chapter of Alpha Sigma Phi must be attained through proper ceremony as prescribed by the Ritual. Every member shall be classified into one of the following categories.
\begin{enumerate}
    \item  \underline{Active Brother} - Once one attains active brother status, he is a member for life. An active  brother is a duly initiated male member who has the following obligations and privileges:
    \begin{itemize}
        \item Must pay all fees associated in due time as determined by the Prudential Board.
        \item Is privileged to be a member of a committee.
        \item Is privileged to be elected to head a committee or hold another chapter office. 
        \item Is privileged to attend, raise issue, and vote at all chapter meetings.
        \item Is privileged to attend all ritualistic events.
        \item Is privileged to attend all social events.
    \end{itemize}
    \item  \underline{New Initiates} - Initiates are newly initiated male members who seek to complete the rest of the initiation process. A man may be initiated as a new initiate as long as ¾ or more of the active brothers vote yes. An initiate is given obligations and privileges within the chapter by the New Member Educator. These privileges shall not include voting in house meetings, unless deemed appropriate as stated in Article 5, or knowledge of the Ritual. Any member of this status can only have their membership removed by a vote of all active brothers to expel them. Furthermore, an initiate can only terminate their status if they choose to expel themselves.
    \item \underline{Alumnus} - An alumnus is a brother who has completed his undergraduate term at RPI or has otherwise left the school. A brother who has left the school without graduating cannot attain this status without approval of the Prudential Board. \textbf{An alumnus is barred from voting in chapter meetings under any circumstances.} Otherwise, an alumnus has the following obligations and privileges:
    \begin{itemize}
        \item Must pay all debts owed to the chapter in that were accrued as an active member.
        \item Is privileged to attend and raise issue at chapter meetings.
        \item Is privileged to attend all ritualistic events.
        \item Is privileged to attend all social events.
        \item Is privileged to visit the chapter at any time.
    \end{itemize}
    The bulleted obligations and privileges above may be modified or taken away from any individual alumnus by a \textbf{three-fifths majority vote by the active brothers}.
    \item  \underline{Little Sister} - A little sister is an honorary position given to a female individual who has garnered a vote from \textbf{three fourths of the active brothers}. They are subject to no obligations from any active brothers or the organization itself and are to be treated with friendship and compassion. \textbf{A little sister is never regarded to be a part of the active chapter and like alumni, are barred from voting during chapter 26 is active right meetings under any circumstances.} A little sister has the following privileges:
    \begin{itemize}
        \item Is privileged to attend and raise issue at chapter meetings.
        \item Is privileged to attend all social events.
        \item Is privileged to visit the chapter at any time except during esoteric events.
    \end{itemize}
    The bulleted privileges above may be modified or taken away from any individual little sister by a \textbf{three-fourths majority vote by the active brothers}.
    \item \underline{Inactive Brother} - An inactive brother holds no obligations and privileges that an ordinary active brother can hold. However he is mandated to pay the required insurance costs that are handed down to us by National Headquarters and any other fines that prude board deems necessary. For all practical purposes, an inactive person is not considered a brother.
\end{enumerate}

\section*{Article 3 - Officer Positions and their Powers}
\hspace*{1cm}The following offices shall be filled at all times by active members of the brotherhood. Any responsibility not listed below shall by default be handled by the president of the chapter until a \textbf{three-fourths majority of the active brothers} officially delegate the said responsibility to another of the officers listed below.\\

\\\underline{HSP (President)} - The President is responsible for and is the official representative of the chapter. He bears the ultimate responsibility of the chapter. He presides over all Chapter Meetings, and his signature is binding for the brotherhood. He is the chief executive officer and responsible for enforcing the constitution, bylaws, and manners of our ritual. He bears a yearlong term of office. His further powers, obligations and privileges are listed below:
\begin{itemize}
    \item Presides over all chapter meetings.
    \item Is the primary owner of the chapter’s bank account along with the HE (Treasurer).
    \item Ensures the timely completion of the responsibilities of all officers
    \item Meets and negotiates with all outside officials (National HQ, RPI, Police, etc...)
    \item Is able to bar any individual, brother or otherwise, from entering property owned by the chapter. This power is to be used at the extreme discretion of the President for the purpose of maintaining the safety and security of the chapter. This ban may only be lifted by the President himself or a \textbf{simple majority of the active brothers.}
    \item The power to declare a temporary time of emergency, and in that time the President may assume any and all powers and responsibilities he deems necessary for the good of the chapter. 
    \item Is a member of the Prudential Board.
    \item Is a member of the Housing Corporation.
\end{itemize}

\underline{HJP (Vice President)} - The Vice President shall have the authority and responsibility of the HSP in his absence. Upon the untimely departure of the HSP from the university or an impeachment of the HSP, the HJP shall fully assume the position of HSP and appoint a new HJP. The Vice President is the chair and secretary of the Prudential Board, is an ex officio member of every committee and is thus responsible for their efficient operation, along with all other internal workings of the chapter. He is responsible for resolving all internal disputes arising among the brothers of the chapter. He bears a yearlong term of office. His further powers, obligations and privileges are listed below:
\begin{itemize}
    \item Presides over Prudential Board meetings.
    \item Ensures the timely completion of the responsibilities of all committee chairs.
    \item Keeps regular correspondence with the Interfraternity Council of RPI.
    \item Assumes all powers and responsibilities of the HSP in his absence.
\end{itemize}

\underline{HE (Treasurer)} - The Treasurer is responsible for and has the authority to collect all brotherhood fees, pay brotherhood debts and invoices, and keep faultless records of all financial transactions. His signature is financially binding for the chapter. He bears a yearlong term of office. His further obligations, powers, and privileges are listed below:
\begin{itemize}
    \item Is the primary owner of the chapter’s bank account, along with the HSP (President).
    \item Ensures payment of all house bills according to a Prudential Board budget.
    \item Ensures payment of all local, state, and federal taxes according to a Prudential Board budget.
    \item Is a member of the Prudential Board.
    \item Is a member of the Housing Corporation.
\end{itemize}

\underline{HS (Secretary)} - The secretary is responsible for keeping an accurate, complete, and impartial account of the proceedings of House Meetings, hereafter referred to as minutes, maintaining a file of all minutes of meetings, and keeping this file accessible to all brothers, keeping custody of a current copy of the constitution and bylaws, and keeping brothers informed of events pertaining to our brotherhood. He must also keep track of the brothers looking to speak during House Meeting discussions to keep order. He bears a semester long term of office.\\

\underline{HC (Sergeant-At-Arms)} - The sergeant-at-arms is the brotherhood chaplain and custodian of all fraternity property pertaining to our ceremonies, paraphernalia, and archived records. He is responsible for the orderly conduct of House Meetings and ceremonies, as well as for protecting the secrets of our brotherhood. He is also an assistant of the HM. He bears a semester long term of office. His further powers, obligations and privileges are listed below:
\begin{itemize}
    \item Will be the Judicial Board delegate for the Interfraternity Council of RPI.
\end{itemize}

\underline{HA (Philanthropy Director)} - The Philanthropy Director is tasked with organizing community service events, coordinating fundraising activities for charitable purposes, and guiding Initiates through their Initiate Class Philanthropy Project.\\

\underline{HP (Alumni Director)} - The alumni director is responsible for all interactions between unaffiliated alumni and the active brotherhood, keeping contact information for alumni, and planning and coordinating alumni activities. He bears a semester long term of office.\\

\underline{HM (Marshall)} - The Marshall is responsible for the ritual. He should have an extensive knowledge of our ritual, and should act in a manner befitting it as an example for the rest of the brotherhood. He is responsible for administering all formal ritualistic events and formal segments of House Meetings. He is also responsible for keeping an accurate, complete, and impartial account of the proceedings of Prudential Board Meetings. He bears a semester long term of office. \\

\underline{HD (New Member Educator)} - The HD is responsible for devising and administering brotherhood education for both initiates and active brothers. He ensures through an initiate   period that initiates are well enough versed in the history, activities, and procedure of the chapter in order to become fully initiated members. He also ensures the active brothers remain contributing members of the organization through lecture and brotherhood development sessions. He bears a semester long term of office.\\

\underline{HZ (Recruitment Director)} - The HZ is responsible for creating, maintaining, and organizing all efforts to recruit new members into the chapter. He holds events to attract interest in the chapter and ensures all members are able and respectable enough to present themselves to and recruit potential candidates.The HZ shall preside over the Recruitment Committee. He bears a semester long term of office.\\

\underline{HCS(Corresponding Secretary)} - The HCS is responsible for maintaining and creating the house's social media. Also Correspondence of activities between other fraternities and sororities. Lastly, serves as one of the members on the rush committee.He bears a semester long term of office.
\section*{Article 4 - Prudential Board}
\hspace*{1cm}The Prudential Board is the legislative, judicial, and financial branch of the chapter. Its membership consists of the HSP, HJP, HD, HM, and HE. This board is responsible for and has the final decision over all financial expenditures. All five members must be present for quorum to be established. However, if a maximum of one prudential office is not filled, the board may meet with four members to conduct business. Under normal circumstances, the HJP is to only use his vote as a tiebreaker. When an office is not filled and the board meets with four members, the HJP may vote or abstain as he wishes. In either case, a matter may only come to pass in the Prudential Board with a minimum of three votes. In the case of one member holding more than one Prudential Office, said member shall only have one vote and a Prude-at-Large shall be elected by the active brothers to ensure there are five individual members on the board. The Prude-at-Large shall remain in office until the end of the semester during which he was elected. This board also has the following powers and obligations:
\begin{itemize}
    \item To allocate an operating budget for the chapter on a regular basis.
    \item The power to allocate funds for various purposes.
    \item The power to levy fines.
    \item The power to impose sanctions.
    \item The power to issue probations.
    \item The power to force a vote for suspensions, expulsions, and for the impeachment of officers at chapter meetings.
    \item The power to act upon all other financial matters related to the chapter.
\end{itemize}
Any action by the Prudential Board may be overturned by a \textbf{simple majority of the active brothers}.
\section*{Article 5 - Chapter Meetings}
\hspace*{1cm}All meetings shall be governed by Robert’s Rules of order whenever it is reasonable to do so, as determined by the HSP or whomever holds the authority of the HSP at the current time. All meetings will be held in an orderly manner. Unless otherwise noted within this constitution or the bylaws of the chapter, all issues pass with a \textbf{simple majority of quorum} in favor, abstentions not counted towards that favor. A meeting must hold quorum (one half of the active brothers plus one) to take place. Once a week the entire brotherhood shall attend a regular chapter meeting, whose date and location shall be determined by the HSP. The chapter meeting is to follow the set order of business as stated:
\begin{enumerate}
    \item \underline{Opening Exercises and Formalities} - The chapter will conduct a meeting closed to all non-active brothers for the purpose of esoteric matters.
    \item \underline{Roll Call} - Each brother will announce his presence after having his name called by the HS for the purpose of establishing quorum.
    \item \underline{Reading of the Previous Minutes} - The HS will read off the previous chapter meeting’s minutes and select portions of other previous meetings that he deems relevant and appropriate to the current meeting.
    \item \underline{Officer Reports} - Each officer will give a brief summary of the activities he was responsible for in the previous week. This may include events, plans made, or correspondence of other parties.
    \item \underline{Committee Head Reports} - Each committee will report their occurrences just as the officers have done.
    \item \underline{Old Business} - Matters that were raised in a previous meeting to be confronted at the current meeting.
    \item \underline{New Business} - Matters not raised previously by the active brothers to be confronted at the current meeting or tabled to the next week’s meeting. The President alone reserves the ability to entertain motions brought up in New Business. Motions the President does not wish to entertain may then be brought up after a simple majority of quorum votes as such.
    \item \underline{Comments and Suggestions for the Good of Society} - Each member is allowed a brief moment to speak their minds on the current state of the brotherhood
\end{enumerate}

Initiates may be granted a vote in non-electoral chapter meetings on a case-by-case basis by a \textbf{simple majority of the active brothers}.

The Chapter Meeting is of utmost importance for the brotherhood, and must be held with exceptions granted for only extreme cases.

For any matter where a voting precedence has not been set previously, the President must set precedence. This precedence becomes constitutional law and must be recorded in the bylaws of the chapter.

\section*{Article 6 - Election Procedure}
\hspace*{1cm}Elections of positions are to be held during the new business segment of a regular Chapter Meeting nearing the end of each semester. Every active brother present maintaining a minimum GPA of 2.7  has the privilege of running for an officer position. However, an exception may be made for a brother who in the previous semester achieved a GPA higher than 2.7 or if the number of active brothers is less than or equal to twenty five. The voting may be separated into however many segments the HSP decides is appropriate. The rules for elections are as follows:
\begin{enumerate}
    \item Election nominations will take place the 10th week of the semester and the elections will be the 11th week.
    \item Brothers must be nominated and seconded at least a week prior to the election of  importance to be considered for a position.
    \item No guests are to be permitted during elections.
    \item Candidates shall be allowed a time of 2 minutes in which to make a speech to the brotherhood. They may then be presented with questions and allowed to answer said questions for 2 minutes. The length of these times is to be enforced by the HC. 
    \item The HS shall count all votes and the HC shall escort candidates away from the meeting before voting commences, and retrieve them when voting is finished.
    \item While the candidates are not present, there may be an open discussion forum about the position in question not exceeding 5 minutes, for HSP, HJP, and HE this will instead not exceed 15 minutes. If the vote fails, another open discussion forum not exceeding 5 minutes will occur, for HSP, HJP, and HE this will instead not exceed 15 minutes. 
    \item A secret ballot may be used when requested by the President of the chapter.
    \item A primary may be held if deemed necessary by the President to limit the number of candidates before proceeding with the final vote. Candidates eliminated by the primary may be allowed to vote in the final election.
    \item For all officer positions, not only must quorum be present for the vote to take place, but a \textbf{simple majority of the present active brothers} must be in favor of a single candidate for the election to be deemed legitimate.
    \item For all other positions, the winner shall be determined by the candidate with the most “For” votes, so long as the abstentions do not total more than one-third of the votes cast.
    \item In the case where there is only one candidate, a motion may be made to elect them to the position by a unanimous vote of all active brothers minus two, for officers, or all present brothers minus two, for other positions.
\end{enumerate}
\section*{Article 7 - Suspension, Expulsion, and Impeachment}
\hspace*{1cm}Upon recommendation by a Prudential Board vote, the brotherhood is forced to take action upon the questions of suspension, expulsion, or removal of office of a brother. Each question may only pass by a \textbf{three-fourths majority of the active brothers}. Suspension, expulsion, or removal of office takes effect immediately upon passing.\\ 

Suspensions and expulsions are to be carried out only on brothers who have committed an egregiously detrimental act against the brotherhood.\\

In the case of an impeachment, elections for the position in question will take place at the next chapter meeting. The duties and responsibilities of the impeached position shall be assumed by the President or, if the President deems appropriate, the Vice President.

\section*{Article 8 – Allegiances}
\hspace*{1cm}All members of the Beta Psi chapter of Alpha Sigma Phi hold allegiance to the 
national Alpha Sigma Phi fraternity and the ideals first put forth by its founding brothers.

\section*{Article 9 – Dissolution}
\hspace*{1cm}In the event that the Beta Psi chapter of Alpha Sigma Phi for any reason cannot continue its operations, the brothers shall hold a vote to dissolve the chapter. With the approval of \textbf{three fourths of the active brothers}, the chapter will cease operations immediately. If the vote to dissolve the chapter succeeds, all chapter assets and records will be transferred to the alumni council or, in the event that no chartered alumni council exists, to the national Alpha Sigma Phi fraternity.

\section*{Article 10 - Amendment}
\hspace*{1cm}This constitution may be amended wholly, or in part by, a \textbf{unanimous vote of the members of the prudential board}. These changes may be overturned by \textbf{simple majority vote by the active members of the chapter} before or during the next meeting.

\section*{Article 11 – Ratification}
\hspace*{1cm}This constitution will be established by a \textbf{unanimous vote} by the voting members of the chapter.

\section*{Article 12 - House Residency}
\hspace*{1cm}The house must have a minimum of one brother living in the house for every room available at the start of the academic year. Rent is to be determined by the Prudential Board. An exemption may be made for those who are a Resident Advisor or Resident Director in an on campus housing location. Residents shall be subject to the Risk Manager’s standards of cleanliness and safety. The President must live in and has supreme precedence in room selection. Brothers who are initiated in the spring semester are not subject to this in the immediate upcoming academic year. A housing lottery shall take place during the 12th week of the semester to determine occupancy and shall be conducted as such:
\begin{itemize}
    \item Living situation shall be decided by roster number. Lower roster number has precedence. 
    \item A choice can be made between living in or living out until there is an equal number of open rooms and brothers whose living situation is undecided.
    \item Doubling of a room may occur if the two occupants mutually agree. Brothers may choose to squat in a room for the upcoming academic year.
    \item Brothers who live-in must be on house meal plan.
    \item Any inactive brother or alumnus wishing to stay in the house with a room assignment must:
    \begin{itemize}
        \item be actively looking for alternative living arrangements
        \item pay rent for each month in advance
        \item pay a \$100 deposit at the beginning of their stay
        \item participate in details
        \item be subject to any rules and regulation that an active brother living in the house would be subject to
        \item contribute significantly (within reason) as either
        \begin{itemize}
            \item house manager
            \item house chef
            \item any scenario not fitting to the stated clauses above can be approved by the Prudential Board
        \end{itemize}
    \end{itemize}
\end{itemize}

\section*{Article 13 - Extended Live-in Guests}
\begin{itemize}
    \item An extended live-in guest is someone who has no other living situation, and is rooming with a brother
    \item Brothers take priority over any guests (I.E., if the brother and his live-in guest want a double, but two other brothers want the same double, the two brothers have precedence over them)
    \item If there is an opening, all extended live-in guests must be approved by a four fifths majority vote of Prude Board
    \begin{itemize}
        \item The brother housing the extended live-in guest will incur a fine equal to 60% of his current rent (That is, the brother will pay 160% of his total rent per semester)
    \end{itemize}
    \item The extended live-in guest is required to do details
    \item The extended live-in guest must be accompanied by a brother in their room
    \item If the extended live-in guest is not accompanied by a brother in the room, then the guest will be given two weeks' notice to find other living arrangements and will consequently be removed from the house.
\end{itemize}

\section*{Article 14- Chain of Command}
\hspace*{1cm}In the event of the absence of a high ranking officer, the following chain of command is to be followed. In the event of a high risk situation the person highest on this chain will assume the responsibilities and obligations of the HSP as dictated in Article 3. This chain was created by unanimous vote of Prudential Board on February 2nd, 2019.
\begin{enumerate}
    \item HSP
    \item HJP
    \item Risk Manager
    \item HC
    \item HM
    \item HE
    \item HD
    \item HZ
    \item HP
    \item HS
    \item HA
    \item All Members of the Standards Board not specified above. (lowest roster number to highest)
    \item All Members of Prudential Board not specified above
    (lowest roster number to highest)
    \item The Active Brother with the lowest roster number who was not specified above. 
\end{enumerate}
\end{document}
